\documentclass[10pt]{article}
\usepackage[a4paper,margin=24mm]{geometry}
\usepackage{amsmath,amssymb,amsthm,booktabs,array}
\usepackage[T1]{fontenc}
\usepackage{lmodern}
\usepackage[colorlinks=true,allcolors=blue,breaklinks=true]{hyperref}
\usepackage{xurl}
\newtheorem{theorem}{Theorem}
\newtheorem{corollary}[theorem]{Corollary}
\newcommand{\cb}[1]{\binom{2#1}{#1}}
\title{Mathar's recurrence in OEIS A174193\\
\large FORMULA, 14 June 2016: a tableau-sum evaluation\\
extending Brignall--Ekhad--Smith--Vatter (2009)}
\author{Haobo Ma\\\small ChronoAI Pte Ltd, Singapore
\and Wenlin Zhang\\\small National University of Singapore\\\small ChronoAI Pte Ltd, Singapore}
\date{4 November 2026}
\begin{document}
\maketitle
\begin{abstract}
We extend the enumeration of Brignall, Ekhad, Smith and Vatter (2009)
by proving the order-four recurrence conjectured by R.\ J.\ Mathar in
OEIS A174193, FORMULA, 14 June 2016, and evaluating their tableau sum.
A central-binomial closed form reduces the
recurrence to an explicit polynomial identity. The same calculation proves
Mathar's recurrence in the cross-referenced entry A174195 for $n\geq4$,
so both recurrences hold for the corresponding permutation counts.
An exact positive difference proves the inequality preceding Problem~10
of Brignall, Ekhad, Smith and Vatter. As a consequence of this inequality,
equal lexicographic ranks give a length-preserving injection.
\end{abstract}

\section{The series and its coefficients}
In the FORMULA section of A174193, Mathar conjectured a recurrence for the
coefficients of the following series \cite{oeis193}:
\begin{align}
 A(x)&=\frac{P(x)-B(x)\sqrt{1-4x}}{2x^2(1-x)(1-4x)^2}
       =\sum_{n\geq0}a_nx^n,\label{eq:series}\\
 P(x)&=1-8x+13x^2+24x^3-48x^4,\notag\\
 B(x)&=1-6x+x^2+34x^3-26x^4-4x^5.\notag
\end{align}
All series are over $\mathbb Q$, and the square root has constant term one.
Brignall, Ekhad, Smith and Vatter give this generating function as an
empirical candidate at the end of Section~2 of \cite{almost}; their
Proposition~3 and the subsequent hook-length calculation give a tableau
characterization and a finite sum for the actual almost-$321$-avoiding permutation counts.
Theorem~\ref{thm:counting} evaluates that sum and identifies it with
\eqref{eq:series}.
The differential-equation approach to recurrences for algebraic series is
classical \cite[Theorems~1.5 and 2.1]{stanley}; Mathar gives explicit constructions
for generating functions involving polynomial roots \cite{mathar}.

Put
\[
 F(t)=t^4-5t^3+10t^2+14t-8,\qquad C_n=\binom{2n}{n}.
\]
\begin{theorem}\label{thm:main}
There is a unique series $A\in\mathbb Q[[x]]$ satisfying
\eqref{eq:series}, and its coefficients are
\begin{equation}
 a_n=\frac{F(n)C_n}{2(n+1)(n+2)(2n-1)}-1\qquad(n\geq0).\label{eq:closed}
\end{equation}
For every integer $n\geq4$ they satisfy Mathar's recurrence
\begin{align}
0={}&-(n+2)(404489n-1866368)a_n\notag\\
&+(3147815n^2-8680082n-14098754)a_{n-1}\notag\\
&+2(-2336439n^2+652022n+14514849)a_{n-2}\notag\\
&+8(-1044770n^2+10585023n-20370553)a_{n-3}\notag\\
&+8(642982n-2004177)(2n-9)a_{n-4}.\label{eq:recurrence}
\end{align}
\end{theorem}
\begin{proof}
The root $s$ with $s(0)=1$ and $s^2=1-4x$ exists by the formal binomial
series and is unique: the sum of two such roots has constant term two and
is a unit. The first two coefficients of $P-Bs$ vanish. After division by
$x^2$, the remaining denominator $2(1-x)(1-4x)^2$ is a unit. This proves
existence and uniqueness in $\mathbb Q[[x]]$; in particular the displayed
division by $x^2$ does not introduce negative powers.

The identities
\[
 P=(1-3x^2)(1-4x)^2,\qquad
 -B=(1-x)(-1+5x+4x^2-30x^3-4x^4)
\]
give the useful Laurent-series decomposition
\begin{equation}
 A=-\frac1{1-x}+\frac1{2x}+\frac1{2x^2}
 +\left(-\frac1{2x^2}+\frac5{2x}+2-15x-2x^2\right)(1-4x)^{-3/2}.
 \label{eq:decomposition}
\end{equation}
The binomial theorem gives
$[x^j](1-4x)^{-3/2}=(2j+1)C_j=:c_j$ for $j\geq0$.
Set $c_j=0$ for $j<0$. Extracting coefficients from
\eqref{eq:decomposition} yields
\[
 a_n=-1-\tfrac12c_{n+2}+\tfrac52c_{n+1}+2c_n-15c_{n-1}-2c_{n-2}.
\]
The ratios needed to simplify this expression are
\begin{equation}
 \frac{(c_{n-2},c_{n-1},c_n,c_{n+1},c_{n+2})}{C_n}
 =\left(\frac{n(n-1)}{4(2n-1)},\frac n2,2n+1,
 \frac{2(2n+1)(2n+3)}{n+1},
 \frac{4(2n+1)(2n+3)(2n+5)}{(n+1)(n+2)}\right).
 \label{eq:ratios}
\end{equation}
These follow from $C_{n+1}/C_n=2(2n+1)/(n+1)$, including $n=0,1$
under the stated zero convention. Substitution proves \eqref{eq:closed}.

Let $q_i(n)$ be the five coefficients in \eqref{eq:recurrence}, in
increasing lag order. Their sum is the zero polynomial, so the constant
$-1$ in \eqref{eq:closed} contributes zero to the recurrence. For the
remaining term use
\[
 \frac{C_{n-i}}{C_n}=\prod_{j=0}^{i-1}\frac{n-j}{2(2n-2j-1)},
 \qquad 0\leq i\leq4.
\]
After multiplication by
$M(n)=4(n+1)(2n-1)(2n-3)(2n-5)(2n-7)$, the five summands become
\begin{align*}
T_0={}&-2(2n-7)(2n-5)(2n-3)(404489n-1866368)F(n),\\
T_1={}&(2n-7)(2n-5)(3147815n^2-8680082n-14098754)F(n-1),\\
T_2={}&-(n+1)(2n-7)(2336439n^2-652022n-14514849)F(n-2),\\
T_3={}&-2n(n+1)(1044770n^2-10585023n+20370553)F(n-3),\\
T_4={}&n(n-1)(n+1)(642982n-2004177)F(n-4).
\end{align*}
Expansion gives $T_0+T_1+T_2+T_3+T_4=0$: all coefficients in degrees
$0$ through $8$ vanish. Since $M(n)\ne0$ for $n\geq4$, this proves
\eqref{eq:recurrence}. The accompanying \texttt{verify.py} checks this
polynomial identity by integer coefficient arithmetic.
\end{proof}

\begin{corollary}\label{cor:first}
The same coefficients satisfy the first-order relation
\[
 (n+3)F(n)(a_{n+1}+1)
 =2(2n-1)F(n+1)(a_n+1)\qquad(n\geq0).
\]
\end{corollary}
\begin{proof}
Insert \eqref{eq:closed} and the ratio $C_{n+1}/C_n$ above.
This proof uses no division by $F(n)$.
\end{proof}

\section{Evaluating the tableau sum}
Write $\operatorname{Av}_n(\beta)^{+1}$ for the permutations of length $n$
from which at most one entry can be deleted to leave a $\beta$-avoider.
For a number $z$, write $(z)_{\underline r}=z(z-1)\cdots(z-r+1)$,
with $(z)_{\underline0}=1$.
\begin{theorem}\label{thm:counting}
For every $n\geq0$, the number
$h_n=|\operatorname{Av}_n(321)^{+1}|$ equals $a_n$ in
\eqref{eq:closed}. Thus \eqref{eq:series} is its counting generating function.
\end{theorem}
\begin{proof}
Proposition~3 and the hook-length calculation in Section~2 of
\cite{almost} give, for $n\geq2$,
\begin{equation}
 h_n=\frac{C_n}{n+1}+\sum_{k=\lfloor n/2\rfloor}^{n-2}
 \left(\frac{n!(2k-n+2)}{(k+2)k!(n-k)(n-k-2)!}\right)^2.
 \label{eq:hook}
\end{equation}
Put $N=n+2$ and $j=n-k$. The summand before squaring becomes
\[
 \frac{(j-1)(N-j-1)(N-2j)}{N(N-1)}\binom Nj.
\]
Its square is unchanged under $j\mapsto N-j$. The terms $j=1,N-1$
vanish, as does the central term when $N$ is even; the terms $j=0,N$
each equal one. Therefore, writing
$w(j)=(j-1)^2(N-j-1)^2(N-2j)^2$, the sum in \eqref{eq:hook} is
\begin{equation}
 \frac{1}{2N^2(N-1)^2}\sum_{j=0}^N w(j)\binom Nj^2-1.
 \label{eq:half-sum}
\end{equation}
Vandermonde's identity gives, for $0\leq r\leq3$,
\begin{equation}
 \sum_{j=0}^N(j)_{\underline r}(N-j)_{\underline r}\binom Nj^2
 =(N)_{\underline r}^2\binom{2N-2r}{N}.
 \label{eq:moments}
\end{equation}
Indeed the left side is $(N)_{\underline r}^2$ times
$\sum_j\binom{N-r}{j-r}\binom{N-r}{j}$; taking the coefficient of
$x^{N-2r}$ in $(1+x)^{2N-2r}$ proves the identity, with out-of-range
binomial coefficients interpreted as zero.
The polynomial identity
\begin{align*}
 w(j)={}&N^2(N-1)^2-N^2(N-1)j(N-j)\\
 &+(N^2-4N+12)(j)_{\underline2}(N-j)_{\underline2}
       -4(j)_{\underline3}(N-j)_{\underline3}
\end{align*}
reduces \eqref{eq:moments} to four terms. Their sum simplifies to
\begin{equation}
 \sum_{j=0}^N w(j)\binom Nj^2
 =\frac{N^2(N-1)^3(N^3-12N^2+48N-60)}
 {4(2N-5)(2N-3)(2N-1)}\binom{2N}{N}.
 \label{eq:weighted}
\end{equation}
This simplification follows by replacing each binomial coefficient in
\eqref{eq:moments} by
$\binom{2N}{N}(N)_{\underline{2r}}/(2N)_{\underline{2r}}$;
all denominators are nonzero for $N\geq4$.
Since $\binom{2N}{N}/C_n=4(2n+1)(2n+3)/((n+1)(n+2))$,
\eqref{eq:hook}--\eqref{eq:weighted} give
\[
 h_n=\frac{C_n}{n+1}
 +\frac{(n^3-6n^2+12n-4)C_n}{2(n+2)(2n-1)}-1=a_n.
\]
For $n=0,1$ both counts and \eqref{eq:closed} are one.
\end{proof}

\section{The cross-referenced series A174195}
A174193 cross-references A174195, whose FORMULA section contains another
recurrence conjectured by Mathar on 14 June 2016 \cite{oeis195}. Write
\begin{align}
 B_*(x)&=\sum_{n\geq0}b_nx^n
 =\frac{R(x)-S(x)(1-4x)^{3/2}}{-2x^2(1-4x)^{3/2}},\label{eq:sibling}\\
 R(x)&=1-5x-6x^2+45x^3-24x^4,\qquad S(x)=1+x-4x^2+x^3.\notag
\end{align}
This is the generating function proved in \cite[Theorem~9]{almost} for
almost-$231$-avoiding permutations.
\begin{theorem}\label{thm:sibling}
Put $G(t)=2t^4-9t^3+19t^2+18t-24$. Then $b_0=b_1=1$, and
\begin{equation}
 b_n=\frac{G(n)C_n}{4(n+1)(n+2)(2n-1)}\qquad(n\geq2).\label{eq:bclosed}
\end{equation}
For $n\geq4$, Mathar's A174195 recurrence holds:
\begin{align}
0={}&2(n+2)(2013n^2-10435n+41550)b_n\notag\\
&+(4026n^3-128025n^2+34955n+59010)b_{n-1}\notag\\
&-2(2n-5)(20130n^2-110855n+93819)b_{n-2}.\label{eq:brec}
\end{align}
The left side is $13998$ at $n=2$ and $57576$ at $n=3$;
thus $n\geq4$ is the exact range among integers $n\geq2$.
\end{theorem}
\begin{proof}
Rewrite \eqref{eq:sibling} as
\[
 B_* =\frac1{2x^2}+\frac1{2x}-2+\frac x2
 +\left(-\frac1{2x^2}+\frac5{2x}+3-\frac{45}2x+12x^2\right)(1-4x)^{-3/2}.
\]
The negative-power coefficients cancel. Coefficient extraction using
\eqref{eq:ratios} gives, for all $n\geq0$,
\[
 b_n=\frac{G(n)C_n}{4(n+1)(n+2)(2n-1)}
       -2\mathbf1_{n=0}+\tfrac12\mathbf1_{n=1}.
\]
For $n\geq4$, all three terms in \eqref{eq:brec} therefore use
\eqref{eq:bclosed}. Multiply the normalized recurrence by
$8(n+1)(2n-3)(2n-1)$. Its numerator is
\begin{align*}
 &4(2n-3)(2013n^2-10435n+41550)G(n)\\
 &\quad +(4026n^3-128025n^2+34955n+59010)G(n-1)\\
 &\quad -(n+1)(20130n^2-110855n+93819)G(n-2)=0.
\end{align*}
All coefficients in degrees $0$ through $7$ vanish. The multiplier is
nonzero in this range, proving the recurrence. At $n=2,3$, substitute
$(b_0,b_1,b_2,b_3)=(1,1,2,6)$ to obtain the stated residuals.
\end{proof}

\begin{corollary}\label{cor:bfirst}
For every $n\geq2$,
\[
 (n+3)G(n)b_{n+1}=2(2n-1)G(n+1)b_n.
\]
\end{corollary}
\begin{proof}
Substitute \eqref{eq:bclosed} and $C_{n+1}/C_n=2(2n+1)/(n+1)$.
\end{proof}

\begin{corollary}\label{cor:difference}
For $n\geq2$,
\[
 b_n-a_n=1+\frac{n-4}{4(2n-1)}\binom{2n}{n}.
\]
In particular
$|\operatorname{Av}_n(231)^{+1}|>|\operatorname{Av}_n(321)^{+1}|$
for $n\geq4$. Both sequences satisfy
$a_n\sim b_n\sim 4^{n-1}\sqrt{n/\pi}$, and
\[
 b_n-a_n\sim\frac{4^n}{8\sqrt{\pi n}},\qquad
 \lim_{n\to\infty}n\frac{b_n-a_n}{b_n}=\frac12.
\]
\end{corollary}
\begin{proof}
Subtract \eqref{eq:closed} from \eqref{eq:bclosed} and use
$G(n)-2F(n)=(n+1)(n+2)(n-4)$. Positivity is immediate for $n\geq4$.
For the asymptotics, Stirling's formula gives $C_n\sim4^n/\sqrt{\pi n}$,
and each rational prefactor in \eqref{eq:closed} and \eqref{eq:bclosed}
is asymptotic to $n/4$. The constant $-1$ is negligible.
The exact difference has central-binomial prefactor tending to $1/8$,
which gives the last two assertions.
\end{proof}

\begin{corollary}\label{cor:injection}
There is a length-preserving injection from
$\operatorname{Av}(321)^{+1}$ to $\operatorname{Av}(231)^{+1}$:
send the $r$th lexicographically ordered member of
$\operatorname{Av}_n(321)^{+1}$ to the $r$th member of
$\operatorname{Av}_n(231)^{+1}$.
\end{corollary}
\begin{proof}
For $n\leq3$ both classes contain every permutation of length $n$.
For $n\geq4$ Corollary~\ref{cor:difference} gives the required cardinality
inequality. Each ordered list is finite and can be constructed by testing
all permutations for avoidance after each single deletion, so the map
is defined and injective at every length. Its definition preserves length.
\end{proof}
This cardinality-induced rank injection does not provide the combinatorial
explanation sought in \cite[Problem~10]{almost}.

\begin{corollary}\label{cor:conjecture11}
Conjecture~11 of \cite{almost} holds for $t=1$: both almost-avoidance
generating functions are rational in $x$ and $\sqrt{1-4x}$.
\end{corollary}
\begin{proof}
Use Theorem~\ref{thm:counting} and \cite[Theorem~9]{almost}.
The case $t=0$ is the classical Catalan case, with generating function
$(1-\sqrt{1-4x})/(2x)$ for both avoidance classes.
\end{proof}

\section{Coverage and verification}
\begin{center}\small
\begin{tabular}{p{.40\linewidth}p{.52\linewidth}}
\toprule
Source statement & Coverage\\\midrule
A174193, FORMULA, 14 June 2016 & Proved: Theorem~\ref{thm:main}, every $n\geq4$, now also for the permutation counts.\\
A174195, FORMULA, 14 June 2016 & Proved for $n\geq4$; nonzero residuals at $n=2,3$ show $n\geq4$ is the exact range.\\
A174195, asymptotic formula, 23 August 2014 & Proved: Corollary~\ref{cor:difference}, including the A174193 analogue and relative gap.\\
\cite{almost}, Section 2, empirical generating function & Proved: Theorem~\ref{thm:counting} evaluates its tableau sum for every $n$.\\
\cite{almost}, Theorem 9 & Prior proof cited: Brignall--Ekhad--Smith--Vatter (2009), the A174195 counting generating function; its closed form is deduced in Theorem~\ref{thm:sibling}.\\
\cite{almost}, inequality before Problem 10 & Proved: the exact count difference is positive for $n\geq4$.\\
\cite{almost}, Problem 10 & Not covered (only the cardinality-induced rank injection is given).\\
\cite{almost}, Conjecture 11 & Proved for $t=1$ (Corollary~\ref{cor:conjecture11}); $t=0$ is the classical Catalan case; $t\geq2$ not covered.\\
\bottomrule
\end{tabular}
\end{center}

The pinned Lean source contains the declaration
\begin{quote}\ttfamily\raggedright
D5.S1.Recurrence.Algebraic.\\
MatharA174193OrderFourRecurrence.result
\end{quote}
It proves existence of the
source series and \eqref{eq:recurrence} for every series satisfying
$s(0)=1$, $s^2=1-4x$, and the cleared equation
$2x^2(1-x)(1-4x)^2A=P-Bs$, for every natural $n\geq4$.
Its exact source is
\href{https://github.com/the-omega-institute/trureturing/blob/0b9787cf2582188d0ab3da380adab3fa26bea39c/D5/S1/Recurrence/Algebraic/MatharA174193OrderFourRecurrence.lean}{\texttt{MatharA174193OrderFourRecurrence.lean}}
at revision \texttt{0b9787cf2582188d0ab3da380adab3fa26bea39c}.
Its imported Euler identity is
$[x^d](x f')=d[x^d]f$ for every $f\in\mathbb Q[[x]]$ and $d\geq0$,
the declaration
\begin{quote}\ttfamily\raggedright
D5.S1.Recurrence.\\A308815DFiniteRecurrence.euler\_coeff
\end{quote}
in the pinned
\href{https://github.com/the-omega-institute/trureturing/blob/0b9787cf2582188d0ab3da380adab3fa26bea39c/D5/S1/Recurrence/A308815DFiniteRecurrence.lean}{\texttt{A308815DFiniteRecurrence.lean}}.
The closed forms, first-order relations, tableau-sum evaluation, A174195
recurrence, comparison, rank injection and $t=0,1$ consequences
are the written proofs given here; they are not assertions of that
\texttt{result} declaration.

The standard-library checker \texttt{verify.py} expands the polynomial
identities in the decompositions, recurrence arguments, coefficient
comparison and weighted Vandermonde calculation.
It independently expands each defining algebraic series through
$x^{200}$, compares all 30 displayed coefficients of each OEIS entry,
checks the two closed forms for $0\leq n\leq200$, the first-order relation
for $0\leq n\leq199$, the A174195 first-order relation for $2\leq n\leq199$,
and both conjectured recurrences for $4\leq n\leq200$.
It compares the tableau sum through $n=200$, enumerates both almost-avoidance
classes by deletion through $n=8$, and checks the rank injection on those lists.
These finite sequence comparisons supplement the universal written proofs.

\begin{thebibliography}{9}
\bibitem{oeis193} N.\ J.\ A. Sloane, entry A174193, \emph{The On-Line
Encyclopedia of Integer Sequences}, 2010; recurrence by R.\ J.\ Mathar,
FORMULA, 14 June 2016. \url{https://oeis.org/A174193}.
\bibitem{almost} R. Brignall, S.\ B. Ekhad, R. Smith and V. Vatter,
Almost avoiding permutations, \emph{Discrete Mathematics} \textbf{309}
(2009), 6626--6631. \href{https://doi.org/10.1016/j.disc.2009.06.027}{doi:10.1016/j.disc.2009.06.027}.
Text: \href{https://arxiv.org/abs/2007.15111v1}{arXiv:2007.15111v1}.
\bibitem{stanley} R.\ P. Stanley, Differentiably finite power series,
\emph{European Journal of Combinatorics} \textbf{1} (1980), 175--188.
\href{https://doi.org/10.1016/S0195-6698(80)80051-5}{doi:10.1016/S0195-6698(80)80051-5}.
\bibitem{mathar} R.\ J. Mathar, P-finite recurrences from generating
functions with roots of polynomials, 2021.
\href{https://arxiv.org/abs/2109.02112v1}{arXiv:2109.02112v1}, Sections 2 and 4.
\bibitem{oeis195} N.\ J.\ A. Sloane, entry A174195, \emph{The On-Line
Encyclopedia of Integer Sequences}, 2010; recurrence by R.\ J.\ Mathar,
FORMULA, 14 June 2016. \url{https://oeis.org/A174195}.
\end{thebibliography}
\noindent\small \copyright\ 2026 Haobo Ma and Wenlin Zhang.
Licensed under Creative Commons Attribution 4.0 International (CC BY 4.0),
\url{https://creativecommons.org/licenses/by/4.0/}.
\end{document}
